\section*{Abstract}
	The T0-Theory (Time-Mass Duality) proposes a geometric reformulation of the foundations of theoretical physics. In simple words: Imagine the universe as a large puzzle whose pieces—from the smallest particles to the vast cosmos—are meant to fit together. The central idea of this work is that \textbf{the natural constants and physical parameters can be traced back to a single dimensionless number} (supplemented by integer structural choices as motivated settings; SI anchors serve for conversion): the universal geometric constant 
	\[
	\xi \approx \frac{4}{3} \times 10^{-4},
	\]
	where $\xi \approx 4/3 \times 10^{-4}$. Imagine $\xi$ as the central parameter—a small number motivated by the basic form of three-dimensional space that links the descriptions of gravity, the speed of light, particle masses, and further quantities.
	
	This collection of over 200 scientific documents develops step by step a physical framework that aims to bring quantum mechanics, relativity, and cosmology together in a common description—based on the principle of absolute time $T_0$ and the intrinsic time-field-mass relationship. In everyday language: The rules of physics are formulated so that time is stable (unlike Einstein's description, in which it is relative), while mass can change; both are connected through this geometric idea. 
	
	The fundamental documents follow a purely geometric path, deriving $\xi$ from the three-dimensional structure of space and tracing further constants back to it, including the fine structure constant 
	\[
	\alpha \approx 1/137,
	\]
	where $\alpha \approx 1/137$, particle masses, and coupling strengths—with $\xi$ as the only parameter (plus the integer settings named in the documents; SI anchors such as $m_e$ serve for conversion). The aim is to need as few arbitrary numbers as possible and to understand the values from a common geometric source. \textbf{[S]} In the cosmic sector, earlier documents postulate a static universe without expansion, as detailed in the \href{https://github.com/jpascher/T0-Time-Mass-Duality/blob/main/2/pdf/140_T0_CMB_Donoghue_Analyse_En.pdf}{CMB document} (verified link exists), in which concepts like dark matter or dark energy would not be needed. The present corpus deliberately leaves the cosmic sector aside, because the Standard Model does not derive it either (P39); these statements are therefore speculative.
	
	This book presents the current state of the T0 Time-Mass Duality Framework and its applications to particle masses, fundamental constants, quantum mechanics, gravity, and cosmology.
	
	The main part of the book consists of a series of core T0 documents. These chapters reflect the current understanding of the theory and its quantitative consequences. Wherever possible, the material has been reorganized and unified to make the structure of the theory as transparent as possible.
	
	The ``Live'' version of the theory is maintained in a public GitHub repository:
	\begin{center}
		\url{https://github.com/jpascher/T0-Time-Mass-Duality}
	\end{center}
	
	The \LaTeX{} sources of the chapters in this book come from this repository. If conceptual or numerical errors are found, they will be corrected there first. This means that the PDF version of the book you are reading is a snapshot of a continuously evolving project. For the most current version of the documents, including new appendices or corrections, the GitHub repository should always be considered the primary reference.
	
	The intention of this compilation is twofold:
	\begin{itemize}
		\item to provide a coherent, readable path through the core ideas and results of the T0-Framework;
		\item to document the historical development of these ideas in the appendix, including false starts, interim formulations, and early adjustments to experimental data.
	\end{itemize}
	
	Readers who are primarily interested in the current formulation of the theory can focus on the core chapters. Readers who are also interested in the considerations and trial-and-error process behind the theory are invited to study the appendix material in parallel.

	\medskip\noindent\textbf{Status markers.} Statements marked \textbf{[S]} are \emph{speculative}: a hint or conjecture that has not yet been derived or numerically checked in the corpus.
	
	\section{The Core Principle: One Central Parameter}
	The basic idea of the T0-Theory can be summarized in one sentence:
	
	\begin{keyresult}
		Central thesis of the T0-Theory: The physical constants—gravitational constant $G$, Planck constant $\hbar$, speed of light $c$, elementary charge $e$, as well as the particle masses and coupling constants—are to be traced back to a single dimensionless number (supplemented by integer structural choices as motivated settings; dimensionful quantities are converted via SI anchors): the universal geometric constant
		\[
		\xi = \frac{4}{3} \times 10^{-4},
		\]
		which emerges from the fundamental three-dimensional space geometry (harmonic fourth 4:3, QFT loop scaling $10^{-4}$).

		From $\xi$ follows the fine structure constant as:
		\[
		\alpha = f_\alpha(\xi) \approx \frac{1}{137.035999177},
		\]
		where $\alpha$ serves as a secondary electromagnetic coupling without primacy.
	\end{keyresult}
	
	In everyday language, this means: The question of where the constants come from is traced back to a single, geometrically motivated number.
	
	\section{Foundations of the T0-Theory}
	\subsection{Time-Mass Duality}
	In contrast to standard physics, where time is relative and mass is constant, the T0-Theory postulates:
	\begin{itemize}
		\item \textbf{Absolute Time Measure} $T_0$: Time flows uniformly everywhere in the universe—like a universal clock that ticks the same for everyone, no matter where you are.
		\item \textbf{Variable Mass}: Mass varies with the energy content of the vacuum—imagine mass as flexible, changing depending on the ``hum'' of the empty space around it.
		\item \textbf{Intrinsic Time Field} $\Tfield$: Every particle carries its own time field—each building block of matter has its personal timer that influences its behavior.
	\end{itemize}
	
	The fundamental relationship is:
	\[
	m(x) = \frac{\hbar}{c^2 \Tfield(x)} = m_0 \cdot (1 + \kappa \Phi(x)),
	\]
	where $\kappa$ is traceable back to $\xi$ via geometric scaling. Mathematically, this duality treats time and mass as variables, and is set up so that the framework remains compatible with established mathematical structures while aiming at a common description of physical phenomena. Simply put: Time and mass are treated as mutually dependent quantities; the mathematics stays manageable and connects to familiar ideas.
	
	\subsection{The Parameter $\xi$}
	The central parameter of the theory is:
	\[
	\xi = \frac{4}{3} \times 10^{-4},
	\]
	a purely geometric construct from 3D space that connects quantum mechanics with gravity. This parameter encodes the fundamental coupling between energy and spatial structure, from which the mass and coupling hierarchies are to be derived. It is like the ratio that tells space how to ``scale'' energy—a small number through which the theory seeks to explain why electrons are light and protons heavy.
	
	\section{Derivation of the Natural Constants}
	\subsection{Tracing Back to $\xi$}
	The T0-Theory traces the following quantities back to $\xi$:
	\begin{enumerate}
		\item \textbf{Gravitational Constant}:
		\[
		G = f_G(\xi, m_P, c, \hbar),
		\]
		where all inputs are reducible to $\xi$-scaled geometric units. Gravity thus appears as an expression of spatial geometry, scaled by $\xi$.
		
		\item \textbf{Particle Masses} (Electron, Muon, Tau, Quarks):
		Particle masses follow a universal scaling law analogous to the ordering principles of atomic energy levels, where quantum numbers $(n, l, j)$ dictate hierarchical structures in a manner similar to atomic shells and subshells—imagine particles stacked like floors in a building, each level set by simple rules, similar to how electrons orbit atoms. Thus,
		\[
		\frac{m_e}{m_P} = g(\xi), \quad \frac{m_\mu}{m_e} = h(\xi), \quad \frac{m_\tau}{m_\mu} = k(\xi),
		\]
		via universal scaling laws $\xi_i = \xi \times f(n_i, l_i, j_i)$. The question of why some particles are about 200 times heavier than others is thus meant to be answered by a common pattern.
		
		\item \textbf{Coupling Constants} (Electroweak, Strong, Electromagnetic):
		\[
		\alpha_W = f_W(\xi), \quad \alpha_s = f_s(\xi), \quad \alpha = f_\alpha(\xi).
		\]
		The coupling strengths are thereby traced back to the same geometric origin.
		
		\item \textbf{Cosmological Parameters} \textbf{[S]}:
		Static universe metrics and CMB temperature $T_{\text{CMB}} = f_{\text{CMB}}(\xi)$, with redshift mechanisms derived from time-field variations (see \href{https://github.com/jpascher/T0-Time-Mass-Duality/blob/main/2/pdf/140_T0_CMB_Donoghue_Analyse_En.pdf}{CMB document} for detailed explanation without expansion). This point is speculative; the present corpus deliberately leaves the cosmic sector aside (P39).
	\end{enumerate}
	
	\section{Experimental Predictions}
	The T0-Theory makes testable predictions:
	
	\begin{foundation}
		Concrete Predictions:
		\begin{itemize}
			\item \textbf{Anomalous Magnetic Moment}: $(g-2)_\mu$ calculation from $\xi$—without additional fit parameters.
			\item \textbf{Koide Formula}: Mass relation of leptons via $\xi$-scaling—the mathematics that connects the weights of three particles in a clean loop.
			\item \textbf{Redshift} \textbf{[S]}: Modified interpretation without expansion, controlled by $\xi$—why distant stars appear ``stretched'' without the universe inflating.
			\item \textbf{CMB Anisotropies} \textbf{[S]}: Interpretation through time-field variations rooted in $\xi$—the microwave ``echo'' of the cosmos as a geometric echo.
		\end{itemize}
	\end{foundation}
	
	These statements are testable—the laboratory quantities with present-day experiments, the cosmological points marked [S] only to a limited extent—and everyone, physicists or curious lay readers, is invited to put the theory to the test.
	
	\section{Structure of the Document Collection}
	This collection includes:
	\begin{itemize}
		\item \textbf{Foundations}: Mathematical formulation of time-mass duality under $\xi$-geometry—the basics explained step by step.
		\item \textbf{Quantum Mechanics}: Deterministic interpretation, Bell inequalities—quantum correlations described deterministically and geometrically local.
		\item \textbf{Quantum Field Theory}: Lagrangian formalism in the T0-Framework—fields in a common description.
		\item \textbf{Cosmology} \textbf{[S]}: Static universe, redshift, CMB—a model without expansion, dark matter, or dark energy (speculative; the present corpus leaves the cosmic sector aside, P39).
		\item \textbf{Particle Physics}: Mass spectrum, anomalous moments, Koide formula—an ordering of the particle spectrum.
		\item \textbf{Technical Applications}: Photon chip, RSA cryptography—possible applications of the theory.
		\item \textbf{Experimental Tests}: Verifiable predictions—tangible ways to investigate the ideas.
	\end{itemize}
	
	Note: The documents consistently follow the geometric $\xi$-path, tracing physics back to 3D space principles, with $\alpha$ and other constants appearing as emergent features. We have woven simple language throughout so that non-experts can dive in without getting lost in jargon.
	
	\section{Epistemological Origin: Four Methodological Convictions}

Before the first formula was written, four methodological convictions stood at the beginning --- not hypotheses to be tested, but working prerequisites that were deliberately set.

\subsection{No Fundamental Contradiction Between QM and RT}

The actual starting point was not a formula but an epistemic stance: quantum mechanics and relativity are both good --- too good to be fundamentally irreconcilable. Two theories that describe reality with such precision should—so the working assumption—not have a genuine contradiction. The apparent conflict is a description artefact, not a natural phenomenon. This conviction guided the search for a common geometric structure --- and led to $\tilde{T} \cdot m = 1$ and thereby to $\xipar$.

\subsection{Instantaneity Must Be a Description Artefact}

Bell tests demonstrate quantum correlations with seemingly instantaneous action across arbitrary distances. But relativity has locality as its foundational structure. If both theories are good, then \emph{instantaneity must be an artefact of the description} --- regardless of how convincingly the Bell tests appear to demonstrate otherwise.

The consequence: Bell correlations are real, but they are geometrically structured on $T^4$, not ontologically nonlocal. Entanglement is a topological connection, not action at a distance. The experimental predictions remain identical to quantum mechanics --- only the interpretation changes \cite{Hilbertraum2026}.

\begin{keyresult}[Key sentence from Dok.~230]
Entanglement as a topological connection --- not as a nonlocal correlation. The experimental predictions are the same.
\end{keyresult}

\subsection{The Simple Torus Is a Simplification}

$T^4$ as a flat identification torus is already a strong simplification of reality. Curvature mechanisms on a non-flat $T^4$ can produce local geometric proximity for infinitely long light paths. What appears in the flat projection as spatially separated points can in the curved geometry be genuinely adjacent.

What appears as instantaneity could therefore be geometric proximity on a curved $T^4$ --- not action at a distance, but local action along a long path that is invisible in the flat projection. This is structurally related to the ER-EPR approach \cite{Maldacena2013}, but without its string-theoretic prerequisites.

\subsection{The Thought Itself Is a $T^4$ Process}

The sharpest form of self-referentiality: the founding thought of FFGFT --- that instantaneity must be an artefact --- is itself a material process within $T^4$. It arose in a physical brain processing structures of reality.

This means: information comes first --- but only because matter comes first, making information processing possible. There is no pure thinking before matter and no pure matter without information processing. FFGFT arose from within $T^4$, by a $T^4$ process, about $T^4$ processes. The theory that says ``we always describe from within'' was itself developed from within \cite{Epistemologie2026}.

This self-referentiality is not regarded here as a weakness, but as an open description of one's own starting position. It also explains why ``why $T^4$?'' remains open: any attempt to answer this question from outside is itself an inside process.

\section{Introduction: Vibrations as the Starting Point}
	The foundation of my T0-Theory did not arise from abstract equations, but from practical work in communications engineering, acoustics, and music theory. Long before I could consider empty space as a dynamic field, I was engaged with vibrations in concrete bodies—for example, the accordion reed \cite{ricot2005}. This small, vibrating membrane in an accordion produces sound through resonance in the ``empty'' air space between: Frequency and amplitude interact dually, without the space remaining ``empty.'' That was an important starting point: Here emergence became tangible—vibration (time) and medium (space) create harmony, without singularities.
	
	This unbiasedness—why not see $\epsilon$ and $\mu$ in QM and EM as dual resonators?—later led to the vacuum approach. In natural units ($\hbar = c = 1$), setting $\alpha$ to 1: EM constants then appear geometric, and a starting point emerges for a common description of QM and RT. The warning against ``translation'' ($\epsilon_0 \neq \mu_0$ naively) was crucial—in T0, $\xi$ ``modulates'' both without loss. From acoustics (resonances in cavities) and communications engineering (Fourier dualities time-frequency \cite{stanfordEE261}) came the entry: Empty space as a resonant vacuum, carried by EM constants ($\epsilon_0$, $\mu_0$, $c = 1/\sqrt{\epsilon_0 \mu_0}$). Music theory reinforced it: Harmonies (Pythagorean 3:4:5 tetrahedra) as fractal overtones hinting at tetra networks.
	
	\section{The Vacuum Approach: From Acoustics to Duality}
	From acoustics (resonances in cavities) and communications engineering (Fourier dualities time-frequency \cite{stanfordEE261}) came the entry: Empty space as a resonant vacuum, carried by EM constants ($\epsilon_0$, $\mu_0$, $c = 1/\sqrt{\epsilon_0 \mu_0}$). Music theory reinforced it: Harmonies (Pythagorean 3:4:5 tetrahedra) as fractal overtones hinting at tetra networks.
	
	T0 formalizes it: The duality $T_{\text{field}} \cdot E_{\text{field}} = 1$ connects time (vibration) and energy (mass), with $\xi$ as the geometric seed. In natural units, set $\alpha = 1$: The Coulomb potential $V(r) = -1/r$ becomes purely geometric, the Bohr radius $a_0 = 1/m_e$ coincides with the reduced Compton length. Tetrahedral networks ``cover'' the time field—emergence of charge/mass without point singularities.
	
	The derivation of $\alpha$:
	\begin{equation}
		\alpha = \xi \cdot \left( \frac{E_0}{1\ \mathrm{MeV}} \right)^2, \quad E_0 = 7.400\ \mathrm{MeV},
	\end{equation}
	yields $\approx 1/137$ \cite{codata2022}, corrected by fractal steps $\prod_{n=1}^{137} (1 + \delta_n \cdot \xi \cdot (4/3)^{n-1})$ to CODATA precision. No ``translation trap''—SI conversion via $S_{\mathrm{T0}} = 1.782662 \times 10^{-30}$ kg projects geometry into the measurement world. Setting $\alpha = 1$ in natural units ($\hbar = c = 1$) makes sense: It reduces EM fluctuations to pure resonance, like in the accordion reed \cite{ricot2005}—vacuum as an acoustic medium where $\epsilon_0$ and $\mu_0$ resonate dually, without naive exchange.
	
	This approach was unbiased: If you set $c = 1$, why not $\alpha$? The consequence: Tetrahedral networks emerge naturally to ``cover'' the time field, and fractal iterations (137 steps) stabilize the emergence of charge and mass. The approach holds because physics can be expressed in dimensionless patterns—from the tangible (vibrations) to the abstract (vacuum).
	
	\section{Convergence with Synergetics: Independent Paths}
	Despite a different approach, T0 and Synergetics converge: Bucky Fuller's tetrahedron as the ``minimum structural system'' \cite{fuller1975} (closest-packing spheres) fractions to vector equilibria—similar to how T0's networks ``pack'' the vacuum. The 137-frequency tetrahedron (2,571,216 vectors = 137 $\times$ 9,384 $\times$ 2) mirrors T0's renormalization: Proton-MeV (938.4) as an emergent ratio.
	
	The independence of the two paths is important: From acoustic resonances (accordion reed as vacuum prototype \cite{ricot2005}) to duality, without Fuller—yet it ``clicks'' at $\alpha=1$. Synergetics offers a parallel that complements the intuitively developed T0 approach: Tetra-fractionation stabilizes vortices (charge), 137 steps as spin transformations (tetra $\to$ octa $\to$ icosa). The long-term engagement with vibrations (accordion reed as resonance example) and unbiasedness ($\epsilon_0$ and $\mu_0$ as dual resonators, without naive translation) independently led to vacuum duality.
	
	\begin{table}[htbp]
		\adjustbox{max width=\textwidth}{%
			\begin{tabular}{lll}
				\toprule
				\textbf{Approach} & \textbf{T0 (Vacuum Duality)} & \textbf{Synergetics (Tetra-Fraction)} \\
				\midrule
				Entry & Acoustics/Resonance in empty space & Closest-Packing Spheres \\
				$\alpha$-Derivation & $\xi \cdot (E_0)^2$ (nat. units: $\alpha=1$) & 137-Frequency Vectors \\
				Time Field & Tetra networks cover duality & Morphological Relativity \\
				Emergence & Charge as vortex (finite $U$) & Vector-Tensor Intertransformation \\
				$\epsilon_0/\mu_0$ & Dual Resonators (modulated via $\xi$) & Tensor Forces in Packing \\
				\bottomrule
		\end{tabular}}
		\caption{Convergences: T0 and Synergetics—extended by duality elements}
		\label{tab:konvergenz}
	\end{table}
	
	From the perspective of this document, the convergence is no coincidence: Both reduce to tetrahedral patterns, but T0 from vacuum resonance (accordion reed as prototype \cite{ricot2005}), Synergetics from packing \cite{fuller1975}. Setting $\alpha=1$ in natural units (Coulomb $V(r) = -1/r$, Bohr radius $a_0 = 1/m_e$) shows: It ``makes sense'' because empty space is geometric—$\epsilon_0$ and $\mu_0$ as dual ``modulators,'' without translation traps. The parallels to Synergetics are a classification made by this document, not a statement by Fuller.
	
	\begin{thebibliography}{9}
		\bibitem{Hilbertraum2026}
	J. Pascher.
	\newblock \emph{Hilbert Space Translation and Bell Correlations in FFGFT} (Dok.~230).
	\newblock Zenodo, 2026.
	\newblock \url{https://doi.org/10.5281/zenodo.20474821}

	\bibitem{Maldacena2013}
	J. Maldacena and L. Susskind.
	\newblock Cool horizons for entangled black holes.
	\newblock \emph{Fortschritte der Physik}, 61(9):781--811, 2013.

	\bibitem{Epistemologie2026}
	J. Pascher.
	\newblock \emph{FFGFT --- Epistemological Position} (Dok.~248).
	\newblock Zenodo, 2026.
	\newblock \url{https://doi.org/10.5281/zenodo.20474821}

	\bibitem{fuller1975}
		R. Buckminster Fuller.
		\newblock \emph{Synergetics: Explorations in the Geometry of Thinking}.
		\newblock Macmillan, 1975.
		
		\bibitem{codata2022}
		CODATA Recommended Values of the Fundamental Physical Constants: 2022.
		\newblock NIST, 2022.
		\newblock URL: \url{https://physics.nist.gov/cuu/pdf/wall_2022.pdf}
		(verified link exists).
		
		\bibitem{ricot2005}
		D. Ricot.
		\newblock The example of the accordion reed.
		\newblock \emph{Journal of the Acoustical Society of America}, 117(4):2279, 2005.
		
		\bibitem{stanfordEE261}
		B. Osgood.
		\newblock \emph{EE 261 - The Fourier Transform and its Applications}.
		\newblock Stanford University, 2007.
		\newblock URL: \url{https://see.stanford.edu/materials/lsoftaee261/book-fall-07.pdf}
		(verified link exists).
	\end{thebibliography}
